% GENERATED FILE. Edit canonical lessons, then run npm run generate:books.
% canonical-source-hash: fnv1a64:2d5e108ce06850a4
% canonical-lessons: MR-C521-bheg, MR-C521-tadakne, MR-C521-gond, MR-C521-karvat, MR-C521-limpne

\chapter{Cracks and Repairs}
\label{ch:mr-cracks-and-repairs}
\hlchaptermodality{521}

\begin{chapteropening}
\textbf{By the end of this chapter:} \emph{I can talk about cracks, glue and simple repairs.}

Complete the last lesson of chapter 521: I can talk about cracks, glue and simple repairs.
\end{chapteropening}

\section[\texorpdfstring{\mr{भेग} (bheg)}{भेग (bheg)}]{\mr{भेग} (bheg) — a fissure; a crack in the ground}

\label{lesson:MR-C521-bheg}



\begin{quote}
Before the new one: say the Marathi for a drain, then the Marathi for to get blocked.
\end{quote}

\subsection*{You'll want to know: \mr{भेग}}
\textbf{\mr{भेग}} — \emph{bheg} — \textquotedblleft{}a fissure; a crack in the ground\textquotedblright{}.

\begin{quote}
\mr{उन्हाळ्यात} \mr{शेतात} \mr{मोठ्या} \mr{भेगा} \mr{पडल्या}.
\end{quote}

\begin{quote}
\emph{unhāḷyāt śetāt moṭhyā bhegā paḍlyā.}
\end{quote}

\begin{quote}
In the summer big cracks opened in the field.
\end{quote}

Feminine. Plural: \mr{भेगा}, \emph{bhegā}.

\mr{भेग} \mr{पडणे}, a crack opens.

\begin{cousinweb}
\textbf{\mr{भेग}} (\emph{bheg}) is the word for a split in dry earth or a wall; its origin is not settled.
\end{cousinweb}

\begin{grammarlens}[title={What you've built}]
One more word for describing what you lived and felt.
\end{grammarlens}

\subsection*{Your turn}
\begin{itemize}
\raggedright
  \item \emph{Say it:} \emph{bheg}
  \item \emph{Say it:} \emph{bheg}, once more
  \item \emph{Say it:} \emph{unhāḷyāt śetāt moṭhyā bhegā paḍlyā.} — In the summer big cracks opened in the field.
  \item \emph{Recall:} say the Marathi for a drain, then the Marathi for to get blocked, then say \emph{bheg} again
\end{itemize}

\begin{tcolorbox}[breakable,colback=teal!4,colframe=teal!35!black,title={Before you move on}]
What does \mr{भेग} mean? (\textquotedblleft{}A fissure; a crack in the ground\textquotedblright{}.) Say it once more.
\end{tcolorbox}
\section[\texorpdfstring{\mr{तडकणे} (taḍakṇe)}{तडकणे (taḍakṇe)}]{\mr{तडकणे} (taḍakṇe) — to crack}

\label{lesson:MR-C521-tadakne}



\begin{quote}
Before the new one: say the Marathi for to get blocked, then the Marathi for a fissure.
\end{quote}

\subsection*{You'll want to know: \mr{तडकणे}}
\textbf{\mr{तडकणे}} — \emph{taḍakṇe} — \textquotedblleft{}to crack\textquotedblright{}.

\begin{quote}
\mr{खाली} \mr{पडल्यावर} \mr{आरसा} \mr{तडकला}.
\end{quote}

\begin{quote}
\emph{khālī paḍlyāvar ārsā taḍaklā.}
\end{quote}

\begin{quote}
The mirror cracked when it fell.
\end{quote}

Past (he …): \mr{तो} \mr{आरसा} \mr{तडकला}, \emph{to ārsā taḍaklā}.

\begin{cousinweb}
\textbf{\mr{तडकणे}} (\emph{taḍakṇe}) is inherited: CDIAL lists Marathi \emph{taḍakṇẽ}, \textquotedblleft{}to crack\textquotedblright{}, with the sound word \emph{traṭ}, \textquotedblleft{}crackle\textquotedblright{} (CDIAL 5988).
\end{cousinweb}

\begin{grammarlens}[title={What you've built}]
One more word for describing what you lived and felt.
\end{grammarlens}

\subsection*{Your turn}
\begin{itemize}
\raggedright
  \item \emph{Say it:} \emph{taḍakṇe}
  \item \emph{Say it:} \emph{taḍakṇe}, once more
  \item \emph{Say it:} \emph{khālī paḍlyāvar ārsā taḍaklā.} — The mirror cracked when it fell.
  \item \emph{Recall:} say the Marathi for to get blocked, then the Marathi for a fissure, then say \emph{taḍakṇe} again
\end{itemize}

\begin{tcolorbox}[breakable,colback=teal!4,colframe=teal!35!black,title={Before you move on}]
What does \mr{तडकणे} mean? (\textquotedblleft{}To crack\textquotedblright{}.) Say it once more.
\end{tcolorbox}
\section[\texorpdfstring{\mr{गोंद} (gond)}{गोंद (gond)}]{\mr{गोंद} (gond) — glue; gum}

\label{lesson:MR-C521-gond}



\begin{quote}
Before the new one: say the Marathi for a fissure, then the Marathi for to crack.
\end{quote}

\subsection*{You'll want to know: \mr{गोंद}}
\textbf{\mr{गोंद}} — \emph{gond} — \textquotedblleft{}glue; gum\textquotedblright{}.

\begin{quote}
\mr{मी} \mr{तुटलेलं} \mr{खेळणं} \mr{गोंदाने} \mr{चिकटवलं}.
\end{quote}

\begin{quote}
\emph{mī tuṭlelaṁ kheḷṇaṁ gondāne cikṭavlaṁ.}
\end{quote}

\begin{quote}
I stuck the broken toy together with glue.
\end{quote}

Masculine.

\begin{cousinweb}
\textbf{\mr{गोंद}} (\emph{gond}) is inherited: CDIAL lists Marathi \emph{gõd}, \textquotedblleft{}gum\textquotedblright{}, beside Hindi \emph{gõd}, under Sanskrit \emph{gundra}, the name of a reed (CDIAL 4199). How a reed's name came to mean gum is not explained there.
\end{cousinweb}

\begin{grammarlens}[title={What you've built}]
One more word for describing what you lived and felt.
\end{grammarlens}

\subsection*{Your turn}
\begin{itemize}
\raggedright
  \item \emph{Say it:} \emph{gond}
  \item \emph{Say it:} \emph{gond}, once more
  \item \emph{Say it:} \emph{mī tuṭlelaṁ kheḷṇaṁ gondāne cikṭavlaṁ.} — I stuck the broken toy together with glue.
  \item \emph{Recall:} say the Marathi for a fissure, then the Marathi for to crack, then say \emph{gond} again
\end{itemize}

\begin{tcolorbox}[breakable,colback=teal!4,colframe=teal!35!black,title={Before you move on}]
What does \mr{गोंद} mean? (\textquotedblleft{}Glue; gum\textquotedblright{}.) Say it once more.
\end{tcolorbox}
\section[\texorpdfstring{\mr{करवत} (karvat)}{करवत (karvat)}]{\mr{करवत} (karvat) — a saw}

\label{lesson:MR-C521-karvat}



\begin{quote}
Before the new one: say the Marathi for to crack, then the Marathi for glue.
\end{quote}

\subsection*{You'll want to know: \mr{करवत}}
\textbf{\mr{करवत}} — \emph{karvat} — \textquotedblleft{}a saw\textquotedblright{}.

\begin{quote}
\mr{सुताराने} \mr{करवतीने} \mr{लाकूड} \mr{कापलं}.
\end{quote}

\begin{quote}
\emph{sutārāne karavtīne lākūḍ kāplaṁ.}
\end{quote}

\begin{quote}
The carpenter cut the wood with a saw.
\end{quote}

Feminine.

The oblique is \mr{करवती}-: \mr{करवतीने}, with a saw.

\begin{cousinweb}
\textbf{\mr{करवत}} (\emph{karvat}) is a tadbhava, inherited from Sanskrit \emph{karapattra}, \textquotedblleft{}saw\textquotedblright{} (CDIAL 2795).
\end{cousinweb}

\begin{grammarlens}[title={What you've built}]
One more word for describing what you lived and felt.
\end{grammarlens}

\subsection*{Your turn}
\begin{itemize}
\raggedright
  \item \emph{Say it:} \emph{karvat}
  \item \emph{Say it:} \emph{karvat}, once more
  \item \emph{Say it:} \emph{sutārāne karavtīne lākūḍ kāplaṁ.} — The carpenter cut the wood with a saw.
  \item \emph{Recall:} say the Marathi for to crack, then the Marathi for glue, then say \emph{karvat} again
\end{itemize}

\begin{tcolorbox}[breakable,colback=teal!4,colframe=teal!35!black,title={Before you move on}]
What does \mr{करवत} mean? (\textquotedblleft{}A saw\textquotedblright{}.) Say it once more.
\end{tcolorbox}
\section[\texorpdfstring{\mr{लिंपणे} (limpṇe)}{लिंपणे (limpṇe)}]{\mr{लिंपणे} (limpṇe) — to plaster; to smear over}

\label{lesson:MR-C521-limpne}



\begin{quote}
Before the new one: say the Marathi for glue, then the Marathi for a saw.
\end{quote}

\subsection*{You'll want to know: \mr{लिंपणे}}
\textbf{\mr{लिंपणे}} — \emph{limpṇe} — \textquotedblleft{}to plaster; to smear over\textquotedblright{}.

\begin{quote}
\mr{गवंड्याने} \mr{भिंतीतली} \mr{भेग} \mr{मातीने} \mr{लिंपली}.
\end{quote}

\begin{quote}
\emph{gavaṇḍyāne bhintītlī bheg mātīne limplī.}
\end{quote}

\begin{quote}
The mason plastered the crack in the wall with mud.
\end{quote}

Past, with -\mr{ने} (he …): \mr{त्याने} \mr{भिंत} \mr{लिंपली}, \emph{tyāne bhint limplī}.

\begin{cousinweb}
\textbf{\mr{लिंपणे}} (\emph{limpṇe}) is a tadbhava, inherited from Sanskrit \emph{limpati}, \textquotedblleft{}smears\textquotedblright{} (CDIAL 11066).
\end{cousinweb}

\begin{grammarlens}[title={What you've built}]
One more word for describing what you lived and felt.
\end{grammarlens}

\subsection*{Your turn}
\begin{itemize}
\raggedright
  \item \emph{Say it:} \emph{limpṇe}
  \item \emph{Say it:} \emph{limpṇe}, once more
  \item \emph{Say it:} \emph{gavaṇḍyāne bhintītlī bheg mātīne limplī.} — The mason plastered the crack in the wall with mud.
  \item \emph{Recall:} say the Marathi for glue, then the Marathi for a saw, then say \emph{limpṇe} again
\end{itemize}

\begin{tcolorbox}[breakable,colback=teal!4,colframe=teal!35!black,title={Before you move on}]
What does \mr{लिंपणे} mean? (\textquotedblleft{}To plaster; to smear over\textquotedblright{}.) Say it once more.
\end{tcolorbox}
